\documentclass[../../main.tex]{subfiles}
\begin{document}

\begin{flushright}
\footnotesize\textit{【自動文字起こし・要確認】}
\end{flushright}

{\bf [解]}
1辺1の正方形$ABCD$を1本の折り目で折り，線対称な五角形$P$（二重になる部分）ができるようにする．図のように，折り返しでできる各点・角を定める（$O$は正方形の対角線の交点とし，折り目は$O$を通る）．$B,C$の像を$B',A'$とし，折り目と辺の交点を$E,F,I,H,G$などとする．$0<x\le1/2,\ 0<\theta\le\pi/4$（$x$は折り返し部分の一辺，$\theta$は折り目の傾き）とする．

求める面積$T$，$\triangle EB'F$の面積$S_1$，$\triangle EA'G(\cong\triangle CEF)$の面積$S_2$とおくと
\begin{align}
T=\dfrac{1}{2}-(S_1+S_2) \quad\cdots\text{②}
\end{align}

$OC=OB'$（対称性）から$\triangle OCB'$は二等辺三角形で$\angle OB'C=\angle OCB'$（③-1）．また$\angle OB'F=\angle OCF=\pi/4$（③-2，正方形の対角線の角度）．③-1,③-2から$\angle B'CF=\angle CB'F$となり，$\triangle FCB'$は二等辺三角形で$CF=B'F$．これと$\triangle IB'F\equiv\triangle ECF$（対称性）から，
\begin{align}
S_1=S_2=\dfrac{1}{2}x^2\tan\theta
\end{align}
②に代入して
\begin{align}
T=\dfrac{1}{2}-x^2\tan\theta \quad\cdots\text{④}
\end{align}

また，辺$\overline{AB}$の長さを$2$通りに表して（折り返し部分の幾何関係から）
\begin{align}
1=x\tan\theta+\dfrac{x}{\cos\theta}+x
\end{align}
\begin{align}
\therefore\ x=\dfrac{\cos\theta}{\sin\theta+\cos\theta+1} \quad(C=\cos\theta,\ S=\sin\theta\ \text{として}\ x=\dfrac{C}{S+C+1}) \quad\cdots\text{⑤}
\end{align}

④から，$T$は$P=x^2\tan\theta$が最大の時最小．⑤から
\begin{align}
P=x^2\tan\theta=\dfrac{C^2}{(S+C+1)^2}\cdot\dfrac{S}{C}=\dfrac{SC}{(S+C+1)^2}
\end{align}
$t=C+S$とおくと，$S+C+1=t+1$，$SC=\dfrac{t^2-1}{2}$（$\because(S+C)^2=S^2+2SC+C^2=1+2SC$）だから
\begin{align}
P=\dfrac{(t^2-1)/2}{(t+1)^2}=\dfrac{t-1}{2(t+1)}=\dfrac{1}{2}\left(1-\dfrac{2}{t+1}\right)
\end{align}
これは$t$の単調増加関数．$t=\sin\theta+\cos\theta=\sqrt{2}\sin\left(\theta+\dfrac{\pi}{4}\right)$は，$0<\theta\le\pi/4$で単調増加し$t\in(1,\sqrt{2}]$．よって$P$は$\theta=\pi/4$（$t=\sqrt{2}$）で最大：
\begin{align}
\max P=\dfrac{1}{2}\left(1-\dfrac{2}{\sqrt{2}+1}\right)
\end{align}
したがって
\begin{align}
\min T=\dfrac{1}{2}-\max P=\dfrac{1}{2}\cdot\dfrac{2}{\sqrt{2}+1}=\dfrac{1}{\sqrt{2}+1}=\sqrt{2}-1
\end{align}


\end{document}
